\subsection{The actual Case II endpoint}
\label{subsec:sr-four-final-endpoint}

We now make the last support and finite-weight-zero passages
explicit. All ordinary inertia characters in this subsection
are specialized to those of the original rational datum:
at \(\ell\) they are the prescribed finite characters times
the indicated cyclotomic powers. They are not independent
ordinary weight variables.

\paragraph{The actual point kills the modular ideal.}
Let \(r\) be an actual ordinary point which is either
absolutely irreducible or a nonsplit block extension with
irreducible diagonal blocks.
Section~\ref{subsec:sr-four-good-residual-lattice} gives,
after a finite coefficient extension and one of the two
permitted labels, its actual good residual lattice.
Choose the tame Case II real tower with this label.
Section~\ref{subsec:sr-four-case-II-residual-links} constructs
the finite residual chain and the finitely many fresh
ramification-killing fields.
Include these fields in the fixed abelian base before
using the Ferrero--Washington bound.
Choose the final cyclotomic layer with
\[
                              d\ge30(C+1)^2 .
 \tag{FE1}
\]
The cohomology constant \(C\) is then fixed throughout the
chain, and the seed is constructed on this final layer.

Take one finite coefficient extension containing the
coefficients of the chain, seed and integral point model.
The actual common space of
Section~\ref{subsec:sr-four-common-space-construction}
receives all these points and their polynomial residual
paths. Section~\ref{subsec:sr-case-II-actual-pm} supplies
one ideal \(J\) in its base \(B\), with the five proved
modular-ideal properties.

The seed's local-zero property gives a minimal component
through the seed's residual point on which \(JA\) vanishes.
Proposition~\ref{prop:sr-four-core}, using (FE1), makes it
vanish on every component through that point.
A residual polynomial path has a prime kernel contained in
the maximal ideals of both endpoints.
Choose a minimal component below that kernel.
If all components at the first endpoint contain \(JA\),
this component does, and hence there is one such component
at the second endpoint. The same proposition makes every
component at the second endpoint contain \(JA\).
Induction along the actual finite chain proves this at the
residual point of the original integral lattice.

In the localization of \(A\) at that residual point,
\(JA\) is therefore in every minimal prime, hence in the
nilradical. The map to the original DVR model extends to
this localization, since elements with nonzero residual
value map to units. Its target is a domain, so it kills
\(J\). This proves the actual point factorization; it
does not assume the point was already automorphic.
Conjugation back to the original representation preserves
characteristic polynomials, purity and automorphy.

In characteristic \(\ell\), its ordinary inertial
characters over the real tower are \(1\).
The finite-weight property of \(J\) then kills
\(\mathfrak m_B\).
Lemma~\ref{lem:sr-four-characteristic-ell} proves that the
original rational characteristic polynomials are exactly
the four prescribed residual characters.
It excludes both an absolutely irreducible rank-four point
and a block extension with an irreducible rank-two block.

In characteristic zero the same factorization, at the fixed
original weight, gives the strict spectral bound (H13).
For an irreducible original point,
Proposition~\ref{prop:sr-case-two-tower-irreducible}
gives irreducibility on the actual selected \(G_L\):
the finite abelian quotient and the residual characters
which remain pairwise distinct there are the needed
hypotheses. Thus the actual classical base change is
cuspidal. Propositions~\ref{prop:sr-all-quadratic-automorphy}
and \ref{prop:sr-actual-wit} give the required rational
quadratic-field automorphy and original-level local data.
This is the all-irreducible-point conclusion with its
restriction and descent premises discharged.

\paragraph{The weight-zero curve has the fixed inertia condition.}
We give the commutative-algebra passage used to deduce finite
weight zero, including the flag at the selected point.

Start with the universal normalized framed representation
with the fixed multiplier and killed-inertia relations.
It is a quotient of a finite-variable complete power-series
ring over the coefficient DVR \(\mathcal O\).
For completeness, finite generation follows from the finite
dimension of continuous \(H^1\) for the finite residual
module and fixed finite ramification set.
Choose lifts of a tangent basis and map a power-series ring
to the universal framed ring. Successive approximation and
completeness make this map surjective.
The normalization equations are closed matrix equations.

Form the local universal framed ring \(\mathcal R_\ell\)
for the same residual representation and multiplier.
Over it, use the projective symplectic full flag variety.
Impose preservation by the local group and the equations
that inertia acts on the \(i\)-th graded line by the
prescribed \(\chi_i\).
This is a closed projective incidence scheme.
The infinitely displayed group equations generate a
coherent ideal, so finitely many suffice.

We select its prescribed residual ordering by an
open-and-closed condition, not by an arbitrary formal tube.
Choose finitely many local group elements \(s_j\) separating
the distinct residual local characters; coincident characters
need no separation.
The characteristic polynomial of \(r(s_j)\) has a unique
Hensel factorization into clusters
\[
             f_j(X)=\prod_\lambda f_{j,\lambda}(X),
 \qquad
             \bar f_{j,\lambda}=(X-\lambda)^{m_{j,\lambda}},
\]
with pairwise coprime factors.
Bezout and Cayley--Hamilton give polynomial spectral
projectors \(e_{j,\lambda}(r(s_j))\).
They preserve every flag in the incidence scheme.
Their ranks on each flag quotient are locally constant;
on a line each idempotent has rank \(0\) or \(1\).
Require rank \(1\) exactly for
\(\lambda=\bar\chi_i(s_j)\) on the \(i\)-th line.
These finitely many rank patterns select an open-and-closed
subscheme \(X_\ell\) of the proper incidence scheme.
It is therefore proper over the complete local base.
Equivalently one may complete this already proper
open-and-closed piece along the whole base special fiber.

Put
\[
 I_\ell=\ker\bigl(\mathcal R_\ell
                         \longrightarrow H^0(X_\ell,\mathcal O)\bigr).
\]
Its proper scheme-theoretic image has underlying set equal
to the proper image.
Take the global quotient by the pullback of \(I_\ell\),
and likewise by the prescribed local image ideals at the
other finitely many places; call the resulting ring \(U\).
Using pullbacks of these local ideals avoids any assumption
that a proper scheme-theoretic image commutes with arbitrary
nonflat base change.

Here is why every original component point maps from \(U\).
Its local \(\mathrm{OnComp}(\mathrm{condL},r_0)\) witness
is a complete local domain \(C\), with the local ordinary
family and a specialization to the split representation
over \(\mathcal O\).
The ordinary condition gives a flag at its generic point
after a finite field extension.
Properness makes the union of the finitely many ordering
images cover \(\operatorname{Spec}C\).
Since \(C\) is a domain, some ordering image is all of
\(\operatorname{Spec}C\).
At the split specialization in characteristic zero the
prescribed inertia weights are distinct
\((a+1,a,1,0)\), so their ordered flag is unique.
Thus this dominating ordering is exactly the one selected
above. Its proper pullback dominates \(\operatorname{Spec}C\),
so \(I_\ell\) maps to zero in \(C\), and hence in the
original coefficient ring.
The other local component conditions are treated the same
way. This proves factorization through \(U\).

We use this closed ordinary locus as an enlargement of the
original component condition.
We do not identify ordinary flags with a choice of
irreducible component.
Conversely, a valuation-valued point of \(U\) belongs
locally to the proper image.
It acquires an actual flag after a finite valued-field
extension, and properness extends the generic flag to
the valuation ring. Its ordered residue characters and
exact fixed inertia characters are those just imposed.
The enlargement suffices: the support exclusion above
holds for every point satisfying these \(KPair\)
conditions, not only the original component condition.

Diagonal change of basis by
\(\operatorname{diag}(t_1t_2,t_2,t_1,1)\) preserves all
these incidence equations.
Thus \(U\) has the complete \(\mathbb Z^2\) grading with
entry weight \(w_{ij}\).
Choose finitely many homogeneous generators \(z_j\) of its
maximal ideal modulo \(\mathfrak m_{\mathcal O}\).
They may be chosen as matrix-entry differences from the
prescribed diagonal representation.
Put
\[
                B_0=U_0,\qquad A_0=\bigoplus_w U_w.
\]
As in the actual common-space construction, the finite
balanced-monomial semigroup shows that \(B_0\) is complete
Noetherian and \(A_0\) is of finite type over \(B_0\).

If \(B_0\) were not finite over \(\mathcal O\), then
\(B_0/\mathfrak m_{\mathcal O}B_0\) would have positive
dimension. Otherwise a power of \(\mathfrak m_{B_0}\)
would lie in \(\mathfrak m_{\mathcal O}B_0\), and complete
Nakayama would make \(B_0\) finite over \(\mathcal O\).
Choose a one-dimensional domain quotient in characteristic
\(\ell\), with nonzero maximal ideal, and normalize it.
This gives a map from \(B_0\) to a complete DVR finite over
\(k[[T]]\), with finite residue field.

Here is the required lift from weight zero to the whole
representation.
For a prime \(\mathfrak p\) of \(B_0\), the grading projection
gives
\(\mathfrak pA_0\cap B_0=\mathfrak p\).
The same holds after localization. Thus the generic fiber
of \(A_0\) over the chosen domain is nonzero.
A closed point in this finite-type fiber has a finite
extension of its fraction field.
Choose a valuation on this extension and write
\(a_j=v(z_j)\) for each nonzero generator image.
Every nonconstant balanced monomial in these generators
lies in \(\mathfrak m_{B_0}\), so has positive valuation.

There is a rational functional \(q:\mathbb Z^2\to\mathbb Q\)
such that
\[
                            a_j+q(w_j)>0
                       \quad\text{for every nonzero }z_j. \tag{FE2}
\]
To verify this elementary strict-inequality assertion, its
failure would give nonnegative rational numbers \(c_j\),
not all zero, with
\(\sum c_jw_j=0\) and \(\sum c_ja_j\le0\), by separation
of the rational polyhedral cone.
Clear denominators. The corresponding nonconstant balanced
monomial would have nonpositive valuation, a contradiction.
After a finite ramified extension, conjugating by the
diagonal torus with parameter valuations \(q(1,0),q(0,1)\)
makes every \(z_j\) have positive valuation.
It leaves every balanced element unchanged.
The map from \(A_0\) therefore extends by completion to
\(U\) and to an actual continuous DVR-valued representation.
Properness of the incidence scheme supplies its actual
ordinary flag and the exact fixed-inertia specialization.

At least one balanced matrix-entry monomial at this point
is not algebraic over \(\mathbb F_\ell\).
Indeed the maximal ideal of \(B_0\) is topologically generated
by \(\mathfrak m_{\mathcal O}\) and finitely many
nonconstant balanced monomials in the \(z_j\).
Some such monomial survives the nonconstant domain quotient.
Expand the diagonal differences in these monomials.
If every resulting balanced matrix-entry monomial were
algebraic over the finite field, their polynomial combination
would be a constant with zero residue, hence zero.
This contradicts the chosen surviving element.
Thus the selected actual point has a balanced monomial \(x\)
such that
\[
                          x^{\ell^n}\ne x
                                      \quad(n>0).        \tag{FE3}
\]
The lift just constructed has exactly the fixed inertia
condition used in (TA7), not a variable-weight ordinary
deformation.

\begin{srproposition}\label{prop:sr-four-actual-final-endpoint}
For every nontrivial principal Case II datum with \(a\ge2\),
the weight-zero algebra is finite over the coefficient DVR;
there is no full irreducible two-block endpoint; and every
irreducible characteristic-zero point has the automorphy
and original-level local properties used in the rank-four
count.
These assertions hold uniformly over finite coefficient
extensions, including coefficient prime \(3\).
\end{srproposition}
\begin{proof}
Suppose finite weight zero failed.
The preceding construction gives an actual
characteristic-\(\ell\) point with (FE3).
Apply Lemma~\ref{lem:sr-four-actual-triangular-alternatives}.

In alternative (I), the actual modular-ideal factorization
proved above and
Lemma~\ref{lem:sr-four-characteristic-ell} exclude the point.
In alternative (D), the fixed-inertia diagonal argument
(TA7)--(TA9) makes every balanced monomial algebraic over
\(\mathbb F_\ell\), contradicting (FE3).
In alternative (B), take the actual block-diagonal part
and apply Proposition~\ref{prop:sr-four-actual-block-extension}.
It gives an actual nonsplit block extension with irreducible
rank-two constituents. The same actual factorization and
characteristic-\(\ell\) lemma exclude this extension.
All three alternatives are impossible.
Consequently \(B_0\) is finite over \(\mathcal O\).

For any original admissible deformation, its balanced
matrix-entry algebra is an image of an \(\mathcal O\)-submodule
of \(B_0\). Since \(\mathcal O\) is Noetherian this algebra is
finite over \(\mathcal O\).
This is the full finite-weight-zero assertion for every
admissible coefficient ring, not merely finiteness of the
reduced characteristic-zero spectrum.

Now suppose a characteristic-zero actual point is block
diagonal with irreducible \(A\)-block.
The block-extension proposition produces a nonsplit actual
extension with the same characteristic polynomials.
Its modular-ideal factorization gives (H13), hence the
same spectral inequality for the original block-diagonal
point.

The ordinary block-projector calculation determines its
two determinant characters on inertia at \(\ell\):
their cyclotomic weights are respectively
\[
                         (a+1)+1=a+2,\qquad a+0=a.
 \tag{FE4}
\]
Each determinant differs from the corresponding cyclotomic
power by a finite character.
Here this follows directly from the local conditions:
after a common finite exponent the quotient character
is trivial on inertia at every finite prime.
At \(\ell\) use the prescribed finite inertial characters;
elsewhere use the actual killed-inertia subgroup of finite
index. At each finite coefficient precision, Minkowski
makes that power character globally trivial.
Completeness gives the same identity integrally, so the
quotient character has finite order.

Choose an open subgroup killing those two finite characters
and lying in the subgroup of (H13).
Chebotarev gives an unramified rational Frobenius in it at
a prime \(p>1\), outside the finite exceptional set.
By (FE4), the ratio of the absolute values of its two block
determinants is \(p^2\).
But if their roots are \(\alpha_1,\alpha_2\) and
\(\beta_1,\beta_2\), (H13) gives
\[
 \frac{|\alpha_1\alpha_2|}{|\beta_1\beta_2|}
       =\frac{|\alpha_1|}{|\beta_1|}
          \frac{|\alpha_2|}{|\beta_2|}<p^2,
\]
a contradiction. This proves the no-full-block conclusion
using the strict bound already proved, without replacing
it by an unproved Ramanujan assertion.

Finally, an irreducible characteristic-zero point factors
through the actual modular ideal as above.
Its irreducibility on the selected CM tower and its descent,
original conductor and local compatibility were proved
there. They give exactly the all-irreducible-point assertion.
Every argument permitted finite coefficient extension;
the arithmetic finite-field comparisons and characteristic
polynomials, rather than division by \(4!\), included
\(\ell=3\).
\end{proof}

The real-tower degree, its class-rank constant and the Hilbert
seed's weight occur only in producing support and the ideal
factorization. No estimate in this endpoint bounds a count
by a function of those auxiliary quantities.
The final spectrum is counted only after descent to the
fixed original rational/imaginary-quadratic level, as in
the uniform nontrivial-character count.

The expanded final source joins are
\srcid{Gsp4PromodularityCaseTwo.last\_step},
\srcid{proMod\_of\_pmEx\_of\_ctx2\_of\_kPair\_of\_star},
\srcid{CompactIntegralFields.PtOver},
\srcid{PtPair}, \srcid{NotOver.bl},
\srcid{DiagConst}, and
\srcid{VWinPm.fin0\_no\_end\_all\_irr\_of\_good\_con\_purity\_aut}.
